\begin{abstract}
Accurate multistep prediction of nonlinear controlled dynamics depends on the
forecast-origin state estimate and the transition model. We study the
Complementary-Filtered Residual World Model (\method), which combines a fixed
gain-0.5, mask-gated context update with a diagonal-Gaussian physics-residual
transition. For correlated prediction and measurement errors, we derive the
channel-wise mean-square-error optimal scalar gain. The confirmatory study
trains \ConfirmatoryRuns{} models across five nonlinear systems, three regimes,
and eight seeds. Against the preselected physics-guided neural-network
comparator, cell-balanced 25-step normalized root-mean-square error (NRMSE-25)
improves by \PrimaryImprovement\% (95\% bootstrap interval
[\PrimaryRelativeCILow\%, \PrimaryRelativeCIHigh\%]; exact one-sided test
against a 5\% margin, $p=\PrimaryPValue$). The aggregate hypothesis is
supported, but clean-condition noninferiority fails. In secondary analyses, a
classical gain-0.5 context filter with nominal rollout attains
$0.111\pm0.003$ NRMSE-25 versus $0.119\pm0.004$ for CFR-WM; disabling the
learned residual also gives the lowest mean among four inference scales. A
synchronous two-step context outage increases CFR-WM error by 2.7\%, versus
51.9\% for the comparator. On held-out measured NanoDrone flights, CFR-WM has
higher full-state NRMSE than a matched unfiltered residual model in all eight
condition--horizon comparisons, but lower 50-step position error in all four;
its nominal 90\% coverage at 50 steps is 0.524--0.667. The results attribute
the controlled-benchmark gain primarily to forecast-origin filtering and show
the limitations of a uniform gain and diagonal uncertainty.
\end{abstract}
